Este trabalho apresenta o desenvolvimento de um sistema web de prontuários eletrônicos para a Liga Odontológica de Reabilitação e DTM da Universidade Estadual de Ponta Grossa (LORD-UEPG). A motivação decorre da gestão inadequada de prontuários físicos em tratamentos contínuos, como os de disfunção temporomandibular (DTM) e de reabilitação protética, que resulta na perda de informações essenciais ao acompanhamento longitudinal. O sistema digitaliza a ficha clínica da Liga a partir de um esquema de dados versionado, persistido em modelo híbrido relacional e documental (PostgreSQL com JSONB), com controle de acesso por papéis, fluxo de supervisão docente, trilha de auditoria e pseudonimização dos dados destinados à pesquisa, em conformidade com a Lei Geral de Proteção de Dados. Como diferencial, o sistema incorpora um módulo de inteligência artificial executado localmente, que transcreve a fala do acadêmico e estrutura as informações narrativas da anamnese nos campos da ficha, sem emitir diagnóstico: toda sugestão é revisada e aprovada por um humano. Um módulo analítico em Python calcula indicadores de evolução clínica, como as escalas de dor, a Escala Graduada de Dor Crônica do RDC/TMD e a amplitude de abertura bucal. A avaliação compreendeu a qualidade da estruturação automática em casos clínicos sintéticos, a comparação entre entrada por áudio e por texto, um teste de usabilidade com membros da Liga e a validação do pipeline analítico com uma coorte sintética de efeito conhecido. \todo{inserir os principais resultados numéricos: F1 da extração, taxa de alucinação, WER, escore SUS e tempo de preenchimento em relação ao papel.}

\textbf{Palavras-chave}: Prontuário eletrônico. Odontologia. Disfunção temporomandibular. Inteligência artificial. Proteção de dados.
